\documentclass[11pt,a4paper]{article}
\usepackage[utf8]{inputenc}
\usepackage[T1]{fontenc}
\usepackage{amsmath}
\usepackage{graphicx}
\usepackage{hyperref}
\usepackage{geometry}
\geometry{a4paper, margin=2.5cm}

\title{Summary: Diffusion Models for MNIST Generation}
\author{}
\date{}

\begin{document}

\maketitle

\noindent GitHub repository: \href{https://github.com/jlvi179/Advanced-Deep-Learning}{https://github.com/jlvi179/Advanced-Deep-Learning}

\section*{Key Concept}

This notebook implements a Denoising Diffusion Probabilistic Model (DDPM) to generate handwritten digits (MNIST). Rather than focusing on a basic implementation from scratch, the notebook utilizes the highly optimized \texttt{denoising\_diffusion\_pytorch} library, demonstrating practical tweaks required to stabilize and accelerate diffusion model training.

\section*{Crucial Tweaks and Optimizations}

Training generative diffusion models is notoriously resource-intensive. The following architectural and procedural tweaks were necessary to make training feasible:

\begin{itemize}
    \item \textbf{U-Net Dimensionality Scaling:} To prevent the model from becoming too large to fit in memory or train in a reasonable time, the base feature dimension was restricted to \texttt{dim = 32}, with a specific multiplier scaling of \texttt{dim\_mults = (1, 2, 5)}. This creates a heavily compressed bottleneck suitable for $28 \times 28$ grayscale images (\texttt{channels = 1}).
    
    \item \textbf{Disabling Flash Attention:} By default, modern transformer/U-Net architectures attempt to use Flash Attention for speed. However, this often leads to compatibility issues on systems without specific CUDA setups or dedicated GPUs. Explicitly setting \texttt{flash\_attn = False} guarantees the code can run across varied hardware environments (like standard Colab or local CPUs).
    
    \item \textbf{Accelerated DDIM Sampling:} The forward diffusion process uses $T = 1000$ steps to smoothly add noise. During inference, stepping back $1000$ times is exceedingly slow. By setting \texttt{sampling\_timesteps = 250}, the model implicitly utilizes DDIM (Denoising Diffusion Implicit Models) sampling, heavily accelerating generation by skipping intermediate timesteps without significantly degrading image quality.
    
    \item \textbf{Weight Decay via AdamW:} Diffusion models are prone to overfitting to the specific noise patterns in smaller datasets. Using \texttt{AdamW} instead of a standard Adam optimizer introduces decoupled weight decay, which helps regularize the U-Net parameters smoothly.
    
    \item \textbf{Progressive Evaluation Visualization:} Because full training requires hours of compute (often exceeding free Colab tiers), a critical workflow tweak is the inline sampling at regular intervals. Using \texttt{make\_grid}, visual outputs are checked dynamically, providing qualitative feedback on convergence long before the final epoch is reached.
\end{itemize}

\end{document}
